\makeabstract
    {
    Pitch Perception in Simultaneous Stream Segregation
    }
    {
    Walter Clemons\textsuperscript{1,2}, Hari Bharadwaj\textsuperscript{2,3}
    }
    {
    \textsuperscript{1} Department of Statistics, Amherst College, Amherst, MA\\
\textsuperscript{2} Center for the Neural Basis of Cognition, Carnegie Mellon University and University of Pittsburgh, Pittsburgh, PA\\
\textsuperscript{3} Department of Communication Science and Disorders, University of Pittsburgh, Pittsburgh, PA
    }
    {
    Individuals with hearing impairment (HI) experience substantial difficulty understanding speech in noisy environments, a challenge known as the ``cocktail party effect.'' Even with hearing aids, this population shows persistent deficits in speech-in-noise comprehension relative to normal-hearing (NH) listeners, indicating that modern-day approaches to hearing deficits do not sufficiently address auditory perception deficits. Pitch presents itself as a possible area for improving these deficits, as HI listeners seem to have poorer pitch perception, and pitch difference between speech and noise has been shown to enable improved speech-in-noise (SiN) comprehension. A better understanding of auditory perception deficits in HI listeners is necessary for improving sensory aids.
    }
    {
    Participants were recruited via Prolific to virtually complete an online SiN task, administered via the SNAPlab online experimental tool. Each stimulus consisted of randomly selected target speech from the Modified Rhyme Test (MRT) simultaneously played with 4 babble utterances from the Harvard Sentence Corpus. Using Praat, babble speech files were manipulated such that their pitch contours followed that of the target speech, appropriately shifted by 0, 2, or 6 semitones. Participants were instructed to listen for the target speech and select the target word. Task performance scores were evaluated across various participant groups and stimulus signal-to-noise ratios. Additionally, NVIDIA{\textquoteright}s Nemotron Speech Streaming ASR model was fine-tuned to complete the MRT task, and its performance was evaluated on the same experimental stimuli.
    }
    {
    SiN comprehension improved significantly with a pitch difference greater than 0, but a 6 semitone pitch difference caused, on average, a significantly smaller improvement. Masking release was slightly, though insignificantly, higher for NH listeners than for HI listeners. However, HI listeners largely demonstrated hearing thresholds within the healthy range.
    }
    {
    Our current findings suggest that SiB comprehension in HI listeners may benefit less from pitch difference. A lack of participants with severe HI may explain limited differences between the two groups. Irregularities in performance at a 6-semitone pitch difference were observed, suggesting possible degradation of audio quality during synthesis. Results from the fine-tuned Nemotron model indicate that pitch perception is related to the physiological separation of distinct sound waves but likely involves higher order cognitive processes.
    }

    \newpage
    
